% =============================================================================
% BMA Health — Burmese (my) language definitions
% =============================================================================
% Report structure
\newcommand{\rptTitle}{ဘန်ကောက်ကျန်းမာရေးစစ်ဆေးမှုဒေတာခွဲခြမ်းစိတ်ဖြာအစီရင်ခံစာ}
\newcommand{\rptSubtitle}{ဘန်ကောက်ကျန်းမာရေးစစ်ဆေးမှုပရောဂျက် --- ဆေးဘက်ဝန်ဆောင်မှုဌာန}
\newcommand{\rptExecSummary}{အမှုဆောင်အကျဉ်းချုပ်}
\newcommand{\rptTotalScreened}{စစ်ဆေးခံသူစုစုပေါင်း}
\newcommand{\rptPersons}{ဦး}
\newcommand{\rptDescriptiveStats}{ဖော်ပြချက်စာရင်းအင်း}
\newcommand{\rptDisease}{ရောဂါ}
\newcommand{\rptDistrict}{ခရိုင်}
\newcommand{\rptZone}{ကျန်းမာရေးဇုန်}
\newcommand{\rptPctAtRisk}{\% အန္တရာယ်ရှိ}
\newcommand{\rptScreened}{စစ်ဆေးခံ}
\newcommand{\rptFactorAnalysis}{အန္တရာယ်အချက်ခွဲခြမ်းစိတ်ဖြာမှု}
\newcommand{\rptInferentialStats}{အကျိုးဆက်စာရင်းအင်း}
\newcommand{\rptTrendAnalysis}{လမ်းကြောင်းခွဲခြမ်းစိတ်ဖြာမှု}
\newcommand{\rptScreeningTests}{ထပ်ဆောင်းစစ်ဆေးမှုရလဒ်}
\newcommand{\rptZoneComparison}{ကျန်းမာရေးဇုန်နှိုင်းယှဉ်ချက်}
\newcommand{\rptPMHealth}{PM2.5 နှင့် ကျန်းမာရေး}
\newcommand{\rptRiskRankings}{ခရိုင်အန္တရာယ်အဆင့်သတ်မှတ်ချက်}
\newcommand{\rptRecommendations}{အကြံပြုချက်များ}
\newcommand{\rptAppendix}{နောက်ဆက်တွဲ}
\newcommand{\rptStatFormulas}{စာရင်းအင်းပုံသေနည်း}
\newcommand{\rptDataSource}{ဒေတာအရင်းအမြစ်}
\newcommand{\rptMethodology}{နည်းစနစ်}
\newcommand{\rptDisclaimer}{စစ်ဆေးခံသူများ၏ဒေတာသာဖြစ်ပြီး လူဦးရေအလုံးစုံကိုယ်စားမပြုပါ}
\newcommand{\rptGenerated}{ထုတ်လုပ်သည့်ရက်}
\newcommand{\rptPage}{စာမျက်နှာ}
\newcommand{\rptOf}{မှ}
% Factor names
\newcommand{\rptFactorSex}{လိင်}
\newcommand{\rptFactorAge}{အသက်အုပ်စု}
\newcommand{\rptFactorOccupation}{အလုပ်အကိုင်}
\newcommand{\rptFactorZone}{ကျန်းမာရေးဇုန်}
\newcommand{\rptFactorBehavior}{ကျန်းမာရေးအပြုအမူ}
\newcommand{\rptSmoking}{ဆေးလိပ်သောက်ခြင်း}
\newcommand{\rptAlcohol}{အရက်သောက်ခြင်း}
\newcommand{\rptExercise}{ကိုယ်လက်လှုပ်ရှားမှု}
% Disease names
\newcommand{\rptDiabetes}{ဆီးချိုရောဂါ}
\newcommand{\rptHypertension}{သွေးတိုးရောဂါ}
\newcommand{\rptDyslipidemia}{သွေးအဆီများခြင်း}
\newcommand{\rptCardiovascular}{နှလုံးသွေးကြောရောဂါ}
\newcommand{\rptStroke}{လေဖြတ်ရောဂါ}
\newcommand{\rptCKD}{နာတာရှည်ကျောက်ကပ်ရောဂါ}
\newcommand{\rptAnemia}{သွေးအားနည်းရောဂါ}
\newcommand{\rptRespiratory}{နာတာရှည်အသက်ရှူလမ်းကြောင်းရောဂါ}
% Screening test names
\newcommand{\rptEKG}{နှလုံးလျှပ်စစ်ဓာတ်ပုံ}
\newcommand{\rptXray}{ရင်ဘတ်ဓာတ်မှန်}
\newcommand{\rptBloodTest}{သွေးစစ်ဆေးမှု}
\newcommand{\rptRetinal}{မျက်စိအကြည်လွှာစစ်ဆေးမှု}
% Statistics terms
\newcommand{\rptChiSquare}{ခိုင်စကွဲယားစမ်းသပ်မှု}
\newcommand{\rptOddsRatio}{အခြေအနေအချိုး}
\newcommand{\rptANOVA}{ကွဲလွဲမှုခွဲခြမ်းစိတ်ဖြာမှု}
\newcommand{\rptLogisticReg}{လော့ဂျစ်စတစ်ဆုတ်ယုတ်ခွဲခြမ်းစိတ်ဖြာမှု}
\newcommand{\rptPValue}{p-တန်ဖိုး}
\newcommand{\rptSignificant}{စာရင်းအင်းအရအရေးပါ}
\newcommand{\rptNotSignificant}{စာရင်းအင်းအရအရေးမပါ}
\newcommand{\rptConfInterval}{95\% ယုံကြည်မှုကြားကာလ}
\newcommand{\rptMannKendall}{Mann-Kendall စမ်းသပ်မှု}
\newcommand{\rptSensSlope}{Sen ၏ slope}
\newcommand{\rptTrend}{လမ်းကြောင်း}
\newcommand{\rptIncreasing}{တိုးလာ}
\newcommand{\rptDecreasing}{ကျဆင်း}
\newcommand{\rptStable}{တည်ငြိမ်}
\newcommand{\rptSeasonal}{ရာသီအလိုက်}
\newcommand{\rptDrySeason}{ခြောက်သွေ့ရာသီ}
\newcommand{\rptRainySeason}{မိုးရာသီ}
